\documentclass[leqno]{article}
% Shared mark scheme preamble for resources.danieljsmith.org: the pack
% preamble, unchanged, then the scheme's own layer. A scheme's .tex is
% rendered from its YAML (scripts/README.md, Mark schemes), never edited.
\input{\CurrentFilePath/pack_preamble.tex}

% The page: equal side margins, so the table runs the full width of the
% header rule. The footers' right edge follows the text's.
\geometry{left=0.8in, right=0.8in}
\makeatletter
\renewcommand{\djs@footR}[1]{%
  \put(\dimexpr1in+\hoffset+\oddsidemargin+\textwidth\relax,
    -\dimexpr\paperheight-0.32in\relax){\makebox[0pt][r]{#1}}}
\makeatother

\usepackage[export]{adjustbox}
\usepackage{tabularray}
\definecolor{msGrid}{RGB}{176, 172, 164}
\definecolor{msTint}{RGB}{245, 243, 239}

% The header on every page, the front page included, and the booklet label
% under the contents-page title. \djsQuestionPageStyle stays empty.
\newcommand{\djsSchemeHeader}[1]{%
  \djsPackHeader[markscheme]{\textit{Mark scheme} - #1}%
  \renewcommand{\djsPackTitle}{#1}%
  \renewcommand{\djsBookletLabel}{Mark scheme}%
}

% The front page: the contents heading, spaced as tocloft spaces the packs'
% "Questions", then the contents table. A question's name and page number
% both link to its first page.
\newcommand{\msContentsHeading}{%
  \vskip\cftbeforetoctitleskip
  {\noindent\cfttoctitlefont\contentsname\par}%
  \vskip\cftaftertoctitleskip}
\NewTblrEnviron{mscontents}
\SetTblrInner[mscontents]{
  colspec = {Q[wd=26mm, l] X[l] Q[wd=16mm, c] Q[wd=14mm, r]},
  width = \linewidth,
  rowsep = 4pt, colsep = 6pt,
  hlines = {0.4pt, msGrid}, vlines = {0.4pt, msGrid},
  hline{2} = {2.2pt, scarlet},
  row{1} = {font=\small\bfseries},
  row{Z} = {font=\bfseries},
  column{2} = {font=\small},
}
\newcommand{\msContentsQuestion}[1]{\hyperref[q:#1]{Question #1}}
\newcommand{\msContentsPage}[1]{\hyperref[q:#1]{\pageref*{q:#1}}}
\newcommand{\msSep}{\hspace{0.9em}}

% The conventions pages' layout: headings in the documents' own bold, tables
% ruled and padded as the scheme's, roomy lists.
\newcommand{\msHead}[1]{%
  \par\vspace{18pt}%
  {\noindent\bfseries #1\par}%
  \vspace{6pt}}
\NewTblrEnviron{mstblr}
\SetTblrInner[mstblr]{rowsep = 6pt, colsep = 6pt,
  hlines = {0.4pt, msGrid}, vlines = {0.4pt, msGrid}}
\newenvironment{msList}
  {\begin{itemize}[leftmargin=1.2em, itemsep=5pt, topsep=3pt, parsep=0pt]}
  {\end{itemize}}

% A question's first page: the link target and the PDF bookmark, as
% \questionitem gives the packs.
\newcommand{\msQnum}{}
\newcommand{\msQuestion}[1]{%
  \clearpage
  \renewcommand{\msQnum}{#1}%
  \phantomsection\label{q:#1}%
  \addcontentsline{toc}{subsection}{\protect\djsTocEntry{#1}}}

% The scheme's table: Question 16mm, Scheme the rest, Marks 16mm, Notes
% 44mm; every cell ruled in msGrid; the scarlet rule under the head row. A
% long table has no caption or "continued" foot; a continuation page opens
% with a muted line above the repeated head row. The rules across a part's
% rows, the alternative's tint and the total row are the rendered table's.
\NewTblrTheme{ms}{%
  \DefTblrTemplate{firsthead}{default}{}%
  \DefTblrTemplate{middlehead,lasthead}{default}{%
    {\small\itshape\color{djsMuted}Question \msQnum\ continued\par}%
    \vspace{6pt}}%
  \DefTblrTemplate{firstfoot,middlefoot,lastfoot}{default}{}%
}
\NewTblrEnviron{msscheme}
\SetTblrOuter[msscheme]{long, theme=ms, entry=none, label=none}
\SetTblrInner[msscheme]{
  colspec = {Q[wd=16mm, l] X[l] Q[wd=16mm, c] Q[wd=44mm, l]},
  width = \linewidth,
  rowsep = 6pt, colsep = 6pt,
  rowhead = 1,
  vlines = {0.4pt, msGrid},
  hline{1} = {0.4pt, msGrid}, hline{2} = {2.2pt, scarlet},
  hline{Z} = {0.4pt, msGrid},
  column{1} = {font=\bfseries}, column{3} = {font=\bfseries},
  column{4} = {font=\small},
  row{1} = {rowsep=4pt, font=\small\bfseries},
}

% Two marks on a row: each note starts with its code.
\newcommand{\msNote}[2]{\textbf{#1}\enspace #2}
% A full-row figure: the worked solution's own picture, scaled uniformly and
% only down, to the table less the Question column. The picture is drawn
% with the line width it has in the solutions document, at its list depth
% there (the optional argument, 1 by default: a part's working), so one
% sized by \linewidth keeps its size; tabularray sets the cell's \linewidth
% to the cell's measured width, which would shrink it at every pass. The
% widths are measured once, in a solution box at the packs' text width
% (pack_preamble.tex's margins, 0.8in and 1in).
\newlength{\msSpanW}
\makeatletter
\newcommand{\ms@keepLW}[1]{%
  \expandafter\xdef\csname ms@solnLW@#1\endcsname{\the\linewidth}}
\AtBeginDocument{%
  \setlength{\msSpanW}{\dimexpr\textwidth-16mm-28pt\relax}%
  \setbox\z@\vbox{%
    \begin{minipage}{\dimexpr\paperwidth-1.8in\relax}
      \begin{djsSolution}
        \ms@keepLW{0}%
        \begin{enumerate}\item \ms@keepLW{1}%
        \begin{enumerate}\item \ms@keepLW{2}%
        \begin{enumerate}\item \ms@keepLW{3}%
        \end{enumerate}\end{enumerate}\end{enumerate}
      \end{djsSolution}
    \end{minipage}}}
\newcommand{\msFigure}[2][1]{%
  \adjustbox{max width=\msSpanW, max totalheight=0.6\textheight, valign=t}{%
    \hsize=\csname ms@solnLW@#1\endcsname
    \linewidth=\hsize
    #2}}
\makeatother
% A figure mark's feature, in the Scheme column.
\newcommand{\msFeature}[1]{\textit{#1}}
% An alternative's title row.
\newcommand{\msAlternative}[1]{\textit{Alternative:} #1}
% Drafts only: a queried mark's code, the query's why under its note, and a
% flag in its part's place.
\newcommand{\msQueried}[2]{#1~\(\to\) #2?}
\newcommand{\msQuery}[1]{{\itshape\color{djsMuted}Query: #1}}
\newcommand{\msFlag}[2]{\textbf{#1:}\enspace #2}


\djsSchemeHeader{Complex Numbers Basics}
\djsSitePath{Core Pure}{Complex Numbers}

\djsLastUpdated{24 October 2026}
\begin{document}
\djsSetQuestionCount{8}
\djsAddMarks{1}{7}
\djsAddMarks{2}{5}
\djsAddMarks{3}{5}
\djsAddMarks{4}{7}
\djsAddMarks{5}{8}
\djsAddMarks{6}{10}
\djsAddMarks{7}{11}
\djsAddMarks{8}{9}
\vspace*{6pt}
\djsContentsTitle
\vspace{30pt}
\msContentsHeading
\begin{mscontents}{}
Question & Marks by part & Marks & Page \\
\msContentsQuestion{1} & (a)(i)\,1\msSep (a)(ii)\,2\msSep (b)(i)\,3\msSep (b)(ii)\,1 & 7 & \msContentsPage{1} \\
\msContentsQuestion{2} & & 5 & \msContentsPage{2} \\
\msContentsQuestion{3} & (a)\,2\msSep (b)\,3 & 5 & \msContentsPage{3} \\
\msContentsQuestion{4} & (a)\,2\msSep (b)\,2\msSep (c)\,3 & 7 & \msContentsPage{4} \\
\msContentsQuestion{5} & (a)\,2\msSep (b)\,6 & 8 & \msContentsPage{5} \\
\msContentsQuestion{6} & (a)(i)\,2\msSep (a)(ii)\,2\msSep (b)(i)\,3\msSep (b)(ii)\,3 & 10 & \msContentsPage{6} \\
\msContentsQuestion{7} & (a)\,4\msSep (b)\,5\msSep (c)\,2 & 11 & \msContentsPage{7} \\
\msContentsQuestion{8} & (a)\,2\msSep (b)\,3\msSep (c)\,2\msSep (d)\,2 & 9 & \msContentsPage{8} \\
\SetCell[c=2]{r} Total & & 62 & \\
\end{mscontents}
\newpage
% The conventions pages, shared by every pack's mark scheme: one source of
% truth, \input after the front page, with no strand or pack logic. The keys
% on the first side, the rules on the second. The scheme layer defines
% \msHead, msList and mstblr.
\vspace*{6pt}%
{\noindent\cfttoctitlefont Reading the mark scheme\par}%
\vspace{10pt}%
These rules apply throughout, unless a note says otherwise.

\msHead{Marks}
Every mark is worth one, so \textbf{M1 A1} on a row is two marks. Only marks
coded ft follow through.\par
\vspace{10pt}%
\begin{mstblr}{
  colspec = {Q[wd=15mm, l] X[l]},
  width = \linewidth,
  hline{2} = {2.2pt, scarlet},
  column{1} = {font=\bfseries},
  row{1} = {font=\small\bfseries, rowsep=4pt},
}
Code & Meaning \\
M1 & Method: a correct method, applied to this question. A slip doesn't cost it if the method is clear. A formula must be used, not just quoted. \\
dM1 & Dependent method: as M1, but only if the method mark before it was earned, unless its note names another. \\
A1 & Accuracy: a correct answer or step. It needs the method mark it follows. \\
A1ft & As A1, following through from an earlier answer that the note names. \\
A1~AG & The last mark of a part whose answer is printed in the question: complete working, with no errors. \\
B1 & A correct value, fact or statement. It needs no method mark. \\
B1ft & As B1, following through from an earlier answer that the note names. \\
E1 & A written explanation, interpretation or conclusion. Concluding a proof by induction needs every earlier mark. \\
\end{mstblr}

\msHead{Abbreviations}
ft and AG qualify a code in the Marks column. The others appear in the
notes.\par
\vspace{10pt}%
\begin{mstblr}{
  colspec = {Q[wd=17mm, l] X[l] Q[wd=15mm, l] X[l]},
  width = \linewidth,
  column{1} = {font=\bfseries}, column{3} = {font=\bfseries},
}
ft & follow through & AG & answer given in the question \\
cao & correct answer only, with no follow-through & cso & correct solution only: no errors in the part \\
oe & or an equivalent form & awrt & any answer that rounds to this value \\
soi & seen, or implied by later working & isw & ignore working after a correct answer \\
SC & special case: credit outside the main method & their & the student's own earlier value \\
condone & a slip that doesn't cost the mark & [\;\,] & working the mark doesn't need \\
s.f. & significant figures & d.p. & decimal places \\
\end{mstblr}

\newpage
\msHead{Reading the scheme}
\begin{msList}
\item Each row is a line of working, with the marks it earns beside it.
  Where a row carries two marks, each note starts with its code.
\item \textbf{[\textit{n}]} ends each part with its total.
\item A shaded \textit{Alternative} is a second method for the same total.
  Mark an answer by the method it follows.
\item Any other valid method earns the equivalent marks, unless the question
  names a method or says \textit{hence}.
\item A sketch earns one B1 for each feature named below its figure.
\end{msList}

\msHead{Accuracy and answers}
\begin{msList}
\item Unless a question says otherwise, give answers to 3 significant
  figures, and angles in degrees to 1 decimal place. Angles in radians take
  3 significant figures.
\item A final answer with more figures is accepted if it rounds to the
  scheme's value; one with fewer is not. An answer spoiled by rounding too
  early loses its accuracy mark.
\item Where a question asks for an exact answer, a decimal loses the final
  accuracy mark. Otherwise either is accepted.
\item Wrong or missing units cost nothing unless the question asks for them.
\item A correct answer with no working earns full marks, unless the question
  needs the working: show that, prove, hence, a method it names, or an exact
  answer. A wrong answer with no working earns nothing. Calculator answers
  are treated the same way.
\item A correct answer, once reached, isn't spoiled by later working (isw).
\end{msList}

\msHead{Value of $\boldsymbol{g}$}
\begin{msList}
\item Take $g = 9.8\,\mathrm{m\,s^{-2}}$ unless a question states a value,
  which is then the only one accepted. Answers that use $g = 9.8$ may be
  given to 2 or 3 significant figures.
\item Where a question doesn't state $g$, an answer worked with $g = 10$
  throughout earns the same marks, keeping any value printed in the
  question. The worked solutions list the final answers with $g = 10$.
\end{msList}


\msQuestion{1}
\begin{msscheme}{
  hline{3,5-6,8-10,12} = {2-Z}{0.4pt, msGrid},
  hline{4,7,11,13} = {0.4pt, msGrid},
  row{13} = {font=\bfseries},
}
Question & Scheme & Marks & Notes \\
1(a)(i) & {\(z^*=x-iy\)} & B1 & \\
\nopagebreak & & [1] & \\
1(a)(ii) & {\(zz^*=(x+iy)(x-iy)\)\\[4pt] \(\phantom{zz^*}=x^2-i^2y^2\)} & M1 & {Expands \((x+iy)(x-iy)\) and uses \(i^2=-1\). Cross terms may cancel unseen; condone one sign slip.} \\
& {\(\phantom{zz^*}=x^2+y^2\), which is real} & A1~AG & {Correct expansion leading without error to \(x^2+y^2\), with a statement that it is real.} \\
\nopagebreak & & [2] & \\
1(b)(i) & {\(2w^*+12=5w-14i\)} & M1 & {Conjugates the whole equation, including \(14i\), or substitutes \(w=a+bi\) and equates real and imaginary parts. Condone one sign slip.} \\
& {\(2w-5w^*=-12+14i\)\\[4pt] \(-5w+2w^*=-12-14i\)\\[4pt] \(\Rightarrow -21w=-84-42i \Rightarrow w=\ldots\)} & dM1 & {Solves their two equations simultaneously, eliminating \(w^*\) or finding \(a\) and \(b\), to reach \(w=\ldots\). Condone arithmetic slips.} \\
& {\(w=4+2i\)} & A1 & {In the form \(a+bi\) only; \(a=4\), \(b=2\) alone does not score.} \\
\nopagebreak & & [3] & \\
1(b)(ii) & {\(w^2(w^*)^2=(12+16i)(12-16i)=400\)} & B1 & \\
\nopagebreak & & [1] & \\
\nopagebreak \SetCell[c=2]{r} Total for Question 1 & & 7 & \\
\end{msscheme}

\msQuestion{2}
\begin{msscheme}{
  hline{3-5} = {2-Z}{0.4pt, msGrid},
  hline{6} = {0.4pt, msGrid},
  row{6} = {font=\bfseries},
}
Question & Scheme & Marks & Notes \\
2 & {\(4(a+bi+2i) = (1-2i)(a-bi)\)\\[4pt] \(\phantom{4}\Rightarrow 4a+(4b+8)i = (a-2b)+(-2a-b)i\)} & M1 & {Substitutes \(z=a+bi\) and \(z^*=a-bi\) and expands the product using \(i^2=-1\). Condone sign slips.} \\
& {\(4a = a-2b\)\\[4pt] \(4b+8 = -2a-b\)} & M1 A1 & {\msNote{M1}{Equates real and imaginary parts of their expansion to form two equations in \(a\) and \(b\).}\\ \msNote{A1}{Correct pair, oe.}} \\
& {\(3a+2b=0,\ 2a+5b=-8 \Rightarrow a=\ldots,\ b=\ldots\)} & dM1 & {Solves their two linear equations for \(a\) and \(b\) by any method; a calculator may be used.} \\
& {\(z=\dfrac{16}{11}-\dfrac{24}{11}i\)} & A1 & {oe, or awrt \(1.45-2.18i\). Must be given as \(z\) in the form \(a+bi\); \(a\) and \(b\) stated separately do not score.} \\
\nopagebreak \SetCell[c=2]{r} Total for Question 2 & & 5 & \\
\end{msscheme}

\msQuestion{3}
\begin{msscheme}{
  hline{3-4,6-8} = {2-Z}{0.4pt, msGrid},
  hline{5,9} = {0.4pt, msGrid},
  row{9} = {font=\bfseries},
}
Question & Scheme & Marks & Notes \\
3(a) & {(i) \(|z_1z_2| = 2\sqrt{3}\)} & B1 & \\
& {(ii) \(\arg(z_1z_2) = -\frac{2\pi}{3}\)} & B1 & {Must lie in \(-\pi < \theta \leq \pi\).} \\
\nopagebreak & & [2] & \\
3(b) & {\(\arg(w^n) = n\left(\text{"}-\tfrac{2\pi}{3}\text{"}\right) = 2k\pi \Rightarrow n = \ldots\)} & M1 & {Multiplies their \(\arg w\) by \(n\) and sets it equal to a multiple of \(2\pi\), or finds arguments of successive powers until one is. Not for a multiple of \(\pi\) alone.} \\
& {(i) \(n = 3\)} & A1 & \\
& {(ii) \(|w^3| = (\text{"}2\sqrt{3}\text{"})^3 = 24\sqrt{3}\)} & B1ft & {ft their \(|z_1z_2|\) raised to their \(n\).} \\
\nopagebreak & & [3] & \\
\nopagebreak \SetCell[c=2]{r} Total for Question 3 & & 5 & \\
\end{msscheme}

\msQuestion{4}
\begin{msscheme}{
  hline{3-4,6-7,9-11} = {2-Z}{0.4pt, msGrid},
  hline{5,8,12} = {0.4pt, msGrid},
  row{12} = {font=\bfseries},
}
Question & Scheme & Marks & Notes \\
4(a) & {\((4+3i)(-12+5i) = -48+20i-36i+15i^2\)} & M1 & {Expands to four terms and uses \(i^2=-1\). Condone one sign slip.} \\
& {\(-63-16i\)} & A1 & \\
\nopagebreak & & [2] & \\
4(b) & {\(|z_2| = \sqrt{(-12)^2+5^2} = 13\)\\[4pt] \(\arg z_2 = \pi - \arctan\frac{5}{12}\)} & M1 & {Attempts the modulus and an argument placed in the second quadrant, such as \(\pi-\arctan\frac5{12}\). Condone a slip in the modulus.} \\
& {\(z_2 = 13(\cos\theta + i\sin\theta)\)\\[4pt] \(\theta = \pi - \arctan\frac{5}{12}\ [\approx 2.75]\)} & A1 & {Form \(13(\cos\theta+i\sin\theta)\) with \(\theta\) exact oe or awrt 2.75.} \\
\nopagebreak & & [2] & \\
4(c) & {\(\arg z_1 = \arctan\frac34\)\\[4pt] \(\arg z_2 = \pi - \arctan\frac5{12}\)\\[4pt] \(\arg(z_1z_2) = \pi + \arctan\frac{16}{63}\)} & B1 & {All three, exact oe or awrt 0.644, 2.75, 3.39, each in \(0\leq\arg z<2\pi\) and its correct quadrant.} \\
& {\(\tan\left(\arctan\frac34 - \arctan\frac5{12}\right)\)\\[4pt] \(\phantom{\tan} = \dfrac{\frac34-\frac5{12}}{1+\frac34\times\frac5{12}}\)} & M1 & {Applies the compound angle formula for \(\tan\) of a sum or difference of their angles, with their tangents substituted. A decimal comparison alone earns neither this nor the A1.} \\
& {\(\phantom{\tan} = \frac{16}{63}\)\\[4pt] \(\arg z_1+\arg z_2 = \pi + \arctan\frac{16}{63}\)\\[4pt] \(\phantom{\arg z_1+\arg z_2} = \arg(z_1z_2)\)} & A1~AG & {Exact \(\frac{16}{63}\) obtained without error, then a statement that the sum of arguments equals \(\arg(z_1z_2)\). Condone no comment on the range of the angle.} \\
\nopagebreak & & [3] & \\
\nopagebreak \SetCell[c=2]{r} Total for Question 4 & & 7 & \\
\end{msscheme}

\msQuestion{5}
\begin{msscheme}{
  hline{3-4,6-10} = {2-Z}{0.4pt, msGrid},
  hline{5,11} = {0.4pt, msGrid},
  row{11} = {font=\bfseries},
}
Question & Scheme & Marks & Notes \\
5(a) & {\(r = \sqrt{(-3)^2 + 2^2}\)\\[4pt] \(\theta = \pi - \arctan\tfrac23\)} & M1 & {Attempts the modulus by Pythagoras and an argument between \(\frac{\pi}{2}\) and \(\pi\) using \(\arctan\frac23\) or \(\arctan\frac32\).} \\
& {\(z_1 = \sqrt{13}\left(\cos 2.554 + i\sin 2.554\right)\)} & A1 & {\(r = \sqrt{13}\) exactly and \(\theta\) awrt 2.554, in the required form.} \\
\nopagebreak & & [2] & \\
5(b) & {\(z_1z_2 = (-3+2i)(a+bi)\)\\[4pt] \(\phantom{z_1z_2} = (-3a-2b) + (2a-3b)i\)} & M1 A1 & {\msNote{M1}{Expands \((-3+2i)(a+bi)\) using \(i^2 = -1\). Condone one sign slip.}\\ \msNote{A1}{Correct real and imaginary parts, oe.}} \\
& {\(2a - 3b = -(-3a-2b) \Rightarrow a = \ldots\)} & M1 & {Uses \(\arg = -\frac{\pi}{4}\) to set their imaginary part equal to minus their real part, or their ratio equal to \(\tan(-\frac{\pi}{4})\), and finds a relation between \(a\) and \(b\).} \\
& {\(a = -5b\)} & A1 & \\
& {\(\sqrt{(13b)^2 + (-13b)^2} = 26\sqrt{2} \Rightarrow b = \ldots\)} & dM1 & {Substitutes their relation into a modulus equation, for \(z_1z_2\) with \(26\sqrt{2}\) or for \(z_2\) with \(2\sqrt{26}\), and solves for \(b\).} \\
& {\(z_2 = -10 + 2i\)} & A1 & {cao. A0 if \(10 - 2i\) is also given as an answer.} \\
\nopagebreak & & [6] & \\
\nopagebreak \SetCell[c=2]{r} Total for Question 5 & & 8 & \\
\end{msscheme}

\msQuestion{6}
\begin{msscheme}{
  hline{3-4,6-7,9-11,13-15} = {2-Z}{0.4pt, msGrid},
  hline{5,8,12,16} = {0.4pt, msGrid},
  row{16} = {font=\bfseries},
}
Question & Scheme & Marks & Notes \\
6(a)(i) & {\(|w| = \sqrt{(-6\sqrt{3})^2 + (-6)^2}\)} & M1 & {Takes the square root of the sum of the squares of the real and imaginary parts. Condone sign slips.} \\
& {\(|w| = 12\)} & A1 & \\
\nopagebreak & & [2] & \\
6(a)(ii) & {\(\arg w = -\pi + \arctan\left(\frac{6}{6\sqrt{3}}\right)\)} & M1 & {Uses arctan of a ratio of the parts' magnitudes and adjusts for the third quadrant, e.g. \(-\pi + \alpha\) or \(-\frac{\pi}{2} - \beta\).} \\
& {\(\arg w = -\frac{5\pi}{6}\)} & A1 & {\(-\frac{5\pi}{6}\) or awrt \(-2.62\), in the range \(-\pi < \arg w \leq \pi\).} \\
\nopagebreak & & [2] & \\
6(b)(i) & {\(z_1z_2 = 2(a\cos\theta + 3\sin\theta)\)\\[4pt] \(\phantom{z_1z_2 =} + 2(a\sin\theta - 3\cos\theta)i\)\\[4pt] \(a\cos\theta + 3\sin\theta = -3\sqrt{3}\)\\[4pt] \(a\sin\theta - 3\cos\theta = -3\)} & M1 & {Expands \(z_1z_2\) and equates real and imaginary parts with \(w\), giving two equations; condone sign slips. Or uses \(|z_1| = \frac{|w|}{|z_2|}\).} \\
& {\((a\cos\theta + 3\sin\theta)^2 + (a\sin\theta - 3\cos\theta)^2\)\\[4pt] \(\phantom{(} = 27 + 9 \Rightarrow a^2 + 9 = 36 \Rightarrow a = \ldots\)} & dM1 & {Squares and adds their equations to eliminate \(\theta\), or sets \(a^2 + 9\) equal to their \(|z_1|^2\), leading to \(a = \ldots\)} \\
& {\(a = 3\sqrt{3}\)} & A1 & {\(3\sqrt{3}\) or \(\sqrt{27}\). A0 if the negative root is also given as an answer.} \\
\nopagebreak & & [3] & \\
6(b)(ii) & {\(z_2 = \frac{-6\sqrt{3} - 6i}{\text{"}3\sqrt{3}\text{"} - 3i} \times \frac{\text{"}3\sqrt{3}\text{"} + 3i}{\text{"}3\sqrt{3}\text{"} + 3i}\)} & M1 & {Divides \(w\) by \(z_1\) with their \(a\), multiplying numerator and denominator by the conjugate of the denominator. Condone sign slips.} \\
& {\(z_2 = -1 - \sqrt{3}\,i\)} & A1 & {Or \(\cos\theta = -\frac12\) and \(\sin\theta = -\frac{\sqrt{3}}{2}\).} \\
& {\(\theta = -\pi + \arctan\sqrt{3} = -\frac{2\pi}{3}\)} & A1 & {In the range \(-\pi < \theta \leq \pi\).} \\
\nopagebreak & & [3] & \\
\nopagebreak \SetCell[c=2]{r} Total for Question 6 & & 10 & \\
\end{msscheme}

\msQuestion{7}
\begin{msscheme}{
  hline{3-6,8-12,14-15} = {2-Z}{0.4pt, msGrid},
  hline{7,13,16} = {0.4pt, msGrid},
  row{16} = {font=\bfseries},
}
Question & Scheme & Marks & Notes \\
7(a) & {\(zw = (3+i)(2+4i)\)\\[4pt] \(\phantom{zw} = 6+12i+2i+4i^2\)} & M1 & {Expands the product using \(i^2=-1\). Condone one arithmetic slip.} \\
& {\(zw = 2+14i\)} & A1 & \\
& {\(|zw| = \sqrt{2^2+14^2} = \ldots\)\\[4pt] \(\arg(zw) = \arctan\frac{14}{2} = \ldots\)} & M1 & {Finds the modulus and argument of their \(zw\), the argument in the correct quadrant for their \(zw\).} \\
& {\(zw = 10\sqrt2\left(\cos(\arctan 7)\right.\)\\[4pt] \(\phantom{zw = 10\sqrt2}\left.+\,i\sin(\arctan 7)\right)\)} & A1 & {Accept \(\sqrt{200}\) or awrt 14.1, and \(\theta\) awrt 1.43.} \\
\nopagebreak & & [4] & \\
7(b) & {\(\frac1v = \frac{z-w}{zw} = \frac{(3+i)-(2+4i)}{\text{"}2+14i\text{"}}\)} & M1 & {Combines over a common denominator, or finds \(\frac1w\) and \(\frac1z\) in the form \(a+bi\) and subtracts.} \\
& {\(\frac1v = \frac{1-3i}{2+14i}\)} & A1 & {oe, such as \(-\frac{2+i}{10}\).} \\
& {\(v = \frac{(2+14i)(1+3i)}{(1-3i)(1+3i)}\)} & M1 & {Inverts their \(\frac1v\) and multiplies numerator and denominator by the conjugate of the denominator.} \\
& {\(\phantom{v} = \frac{2+6i+14i+42i^2}{1-9i^2}\)} & dM1 & {Expands numerator and denominator using \(i^2=-1\), reaching a real denominator. Condone one slip.} \\
& {\(v = -4+2i\)} & A1 & \\
\nopagebreak & & [5] & \\
7(c) & {\(vv^* = (-4+2i)(-4-2i) = \ldots\)} & M1 & {Multiplies their \(v\) by its conjugate, or uses \(vv^*=a^2+b^2\) with their \(a\) and \(b\).} \\
& {\(vv^* = 20\), so \(k = 20\)} & A1 & \\
\nopagebreak & & [2] & \\
\nopagebreak \SetCell[c=2]{r} Total for Question 7 & & 11 & \\
\end{msscheme}

\msQuestion{8}
\begin{msscheme}{
  hline{3-4,6-8,10-11,13-15} = {2-Z}{0.4pt, msGrid},
  hline{5,9,12,16} = {0.4pt, msGrid},
  row{16} = {font=\bfseries},
}
Question & Scheme & Marks & Notes \\
8(a) & {Real coefficients \(\Rightarrow\) \(3+i\) and \(2+3i\) also roots} & E1 & {States that non-real roots occur in conjugate pairs, so the conjugates of \(z_2\) and \(z_3\) must also be roots.} \\
& {Five distinct roots, but a quartic has at most four} & E1 & {Concludes that five distinct roots, or more than four, is impossible for a degree-four polynomial.} \\
\nopagebreak & & [2] & \\
8(b) & {\(\frac{2-i}{1-3i} \times \frac{1+3i}{1+3i} = \ldots\)} & M1 & {Substitutes and multiplies numerator and denominator by the conjugate of their denominator. Condone sign slips in forming the differences.} \\
& {\(\phantom{\frac{2}{1}} = \frac12 + \frac12i\)} & A1 & {oe, such as \(\frac{5+5i}{10}\).} \\
& {First quadrant \(\Rightarrow \arg = \arctan 1 = \frac{\pi}{4}\)} & A1~AG & {Correct quotient leading without error to \(\frac{\pi}{4}\), with the first quadrant, or both parts positive, stated.} \\
\nopagebreak & & [3] & \\
8(c) & {\(\arg\left(\frac{z_2-z_1}{z_3-z_1}\right) = \arg(2-i) - \arg(1-3i)\)\\[4pt] \(\phantom{\arg} = -\arctan\frac12 - \left(-\arctan 3\right)\)} & M1 & {Uses the argument of a quotient as the difference of arguments, with arctan forms of both arguments placed in the fourth quadrant. Condone one sign slip.} \\
& {\(\frac{\pi}{4} = \arctan 3 - \arctan\frac12\)} & A1~AG & {Correct negative arguments equated to \(\frac{\pi}{4}\) from (b), reaching the printed result without error.} \\
\nopagebreak & & [2] & \\
8(d) & \SetCell[c=3]{c} \msFigure{\begin{tikzpicture}[scale=1.35]
\def\xmin{-1}
\def\xmax{5}
\def\ymin{-4.5}
\def\ymax{2}
\def\xone{1}
\def\yone{0}
\def\xtwo{3}
\def\ytwo{-1}
\def\xthree{2}
\def\ythree{-3}
\pgfmathsetmacro{\yleft}{(\xmin-6)/3}
\pgfmathsetmacro{\yright}{(\xmax-6)/3}
\coordinate (O) at (0,0);
\coordinate (z1) at (\xone,\yone);
\coordinate (z2) at (\xtwo,\ytwo);
\coordinate (z3) at (\xthree,\ythree);
\coordinate (L) at (\xmin,\yleft);
\coordinate (R) at (\xmax,\yright);
\begin{scope}
\clip (\xmin,\ymin) rectangle (\xmax,\ymax);
\fill[gray!30] (L) -- (R) -- (\xmax,\ymax) -- (\xmin,\ymax) -- cycle;
\end{scope}
\draw[thick, black, -{Stealth[length=3mm]}] (\xmin,0) -- (\xmax,0) node[below] {$\mathrm{Re}$};
\draw[thick, black, -{Stealth[length=3mm]}] (0,\ymin) -- (0,\ymax) node[left] {$\mathrm{Im}$};
\node[below left] at (O) {$O$};
\draw[thick, black] (L) -- (R);
\foreach \p/\lab in {z1/z_1,z2/z_2,z3/z_3} {
  \fill[black] (\p) circle (1.5pt);
  \node[above] at (\p) {$\lab$};
}
\end{tikzpicture}} & & \\
& \msFeature{Shape} & B1 & {Straight line through the midpoint of \(z_1\) and \(z_3\), visibly perpendicular to their join. Intercepts need not be marked.} \\
& \msFeature{Shading} & B1 & {Half-plane containing \(z_1\) shaded, on a line that is at least roughly the bisector. Condone a dashed boundary.} \\
\nopagebreak & & [2] & \\
\nopagebreak \SetCell[c=2]{r} Total for Question 8 & & 9 & \\
\end{msscheme}

\end{document}
